%! Piecewise \& Step Functions — Unit 1, Lesson 1.4 — Student Packet Cover Sheet
% Regenerated 2026-09-07 in the course-wide gradual-release shape: three scored rows (Warm-Up,
% Guided Notes, Homework), no activity row, no debrief row. Homework is due the first class after
% two study halls — course policy of record (2026-09-03) — never "next class".
\documentclass[10pt]{article}
\usepackage{algebra2-article}
\usepackage{algebra2-boxes}
\usepackage{ltablex}
\keepXColumns

\begin{document}

% Full-bleed forest banner. The header text is set INSIDE a tikz node anchored to the page so it
% can never drift above the paper edge.
\begin{tikzpicture}[remember picture, overlay]
  \fill[forest] (current page.north west)
    rectangle ([yshift=-1.16in]current page.north east);
  \node[anchor=north, align=center, inner sep=0pt, text width=\paperwidth]
    at ([yshift=-0.28in]current page.north) {%
      {\color{white}\bfseries\LARGE Algebra 2}\\[4pt]
      {\color{forestmid}\bfseries\large Unit 1 \quad Linear Functions}\\[6pt]
      {\color{white}\bfseries\large Lesson 1.4 \quad Piecewise \& Step Functions}%
    };
\end{tikzpicture}

% Push the body clear of the 1.16in banner (page top margin is 0.60in).
\vspace*{0.56in}
\namedateperiod
\vspace{0.12in}

\begin{learningtargetbox}
\small
\begin{enumerate}[leftmargin=1.5em, itemsep=5pt]
  \item \ldots{} read a \textbf{piecewise-defined function} as \emph{one} function with different
        rules on different pieces of its domain, and \textbf{evaluate} it by first choosing the
        piece the input belongs to --- at a \textbf{boundary}, the piece whose inequality
        \emph{includes} the input.
  \item \ldots{} write the absolute value function as a \textbf{two-piece} linear rule, so that
        Lesson 1.3's V is the first piecewise function I already know.
  \item \ldots{} read a piecewise graph using \textbf{open} and \textbf{closed} endpoints, and
        decide whether the pieces \emph{meet} at a boundary or \emph{jump} (a
        \textbf{discontinuity}).
  \item \ldots{} evaluate the \textbf{greatest-integer (step) function} $\lfloor x\rfloor$ by
        rounding \emph{down}, read its staircase graph, and name \textbf{constant} intervals.
\end{enumerate}
\end{learningtargetbox}

\vspace{0.14in}

\begin{tocbox}
\small
\renewcommand{\arraystretch}{1.5}
\begin{tabularx}{\linewidth}{@{} c l X r @{}}
\rowcolor{forestbg}
\textbf{\#} & \textbf{Component} & \textbf{Description} & \textbf{Score} \\
\midrule
% THREE scored rows in packet order. No activity row (there is no group activity), no debrief row
% (a phase, not a component). Four columns: every row needs FOUR cells or the widths collapse.
1 & Warm-Up        & Spiral review: evaluating a rule, the V as two lines, and open vs.\ closed
                     endpoints                                                                 & \blank{1.2cm} \\
2 & Guided Notes   & A streaming service, free for three months: one function with several rules;
                     the V as two pieces; meets or jumps; the greatest-integer staircase --- ending
                     in \emph{Guided Practice} and \emph{Individual Practice}                  & \blank{1.2cm} \\
3 & Homework       & Evaluating with a boundary, a meets-or-jumps pair, reading a graph with a
                     jump, the staircase, a phone-plan model, and one multiple-choice item ---
                     \textbf{scored; due the first class after two study halls}                & \blank{1.2cm} \\
\midrule
  & \multicolumn{2}{r}{\textbf{Total}} & \blank{1.2cm} \\
\end{tabularx}
\end{tocbox}

\vspace{0.10in}

\begin{remindbox}
\small A \textbf{piecewise-defined function} is \emph{one} function that follows different rules
on different parts of its domain. To \textbf{evaluate} it, first find the piece whose condition
the input satisfies, then apply that rule. At a \textbf{boundary}, exactly one piece owns the
input --- the one whose inequality \emph{includes} it ($\le$ or $\ge$), never the strict one, and
never ``whichever the story makes sound right.'' On a graph a \textbf{closed} dot includes its
endpoint and an \textbf{open} dot excludes it; the pieces \emph{meet} at a boundary only when both
rules give the same value there, and otherwise the graph \emph{jumps} --- a
\textbf{discontinuity}. The \textbf{greatest-integer function} $\lfloor x\rfloor$ rounds
\emph{down} (so $\lfloor -2.5\rfloor=-3$); its graph is a staircase, \textbf{constant} on each
interval $[n,n+1)$.
\end{remindbox}

\end{document}
